% Shared preamble of the talk. Loaded by slides.tex (beamer, \decktrue) and
% summary.tex (article, \deckfalse). Standalone: it does not use a paper's main.tex.
\usepackage[utf8]{inputenc}
\usepackage[T1]{fontenc}
\usepackage{lmodern}
\usepackage{amsmath,amssymb,amsfonts}
\usepackage{xcolor}
\usepackage{tikz}
\usepackage{adjustbox}
\usepackage{graphicx}
\usetikzlibrary{arrows.meta,calc,positioning,decorations.markings}

% Which parts are built. \Bonusfalse removes the bonus frames from the deck; \Backupfalse the appendix.
\newif\ifBonus \Bonustrue
\newif\ifBackup \Backuptrue


% ---------- maths macros (add the talk's own here; one symbol per object) ----------
\newcommand{\R}{\mathbb{R}}
\newcommand{\Z}{\mathbb{Z}}
\newcommand{\Q}{\mathbb{Q}}

% ---------- no status tags and no status notes: a cited result names its source in the
% environment header, an unproved statement is an `idea` environment (see SKILL.md).

% ---------- dual-mode helpers ----------
\ifdeck
  \newcommand{\onlydeck}[1]{#1}
  \newcommand{\onlyarticle}[1]{}
  \newcommand{\slide}[2]{\begin{frame}[t]\frametitle{#1}\vspace{0.3em}#2\end{frame}}
  \newcommand{\bslide}[2]{\slide{Bonus: #1}{#2}}
  \newcommand{\aslide}[2]{\slide{Appendix: #1}{#2}}
  \newcommand{\dslide}[2]{\slide{#1}{#2}}
  \newcommand{\slideT}[3]{\slide{#1}{#3}}
  \newcommand{\appref}{the appendix}
  \newcommand{\speak}[1]{\note{#1}}
\else
  \newcommand{\onlydeck}[1]{}
  \newcommand{\onlyarticle}[1]{#1}
  \newcommand{\slide}[2]{\subsection{#1}#2\par}
  \newcommand{\bslide}[2]{\slide{#1}{#2}}
  \newcommand{\aslide}[2]{\slide{#1}{#2}}
  \newcommand{\dslide}[2]{}
  \newcommand{\slideT}[3]{\slide{#2}{#3}}
  \newcommand{\appref}{the bonus section}
  \newcommand{\speak}[1]{}
\fi

% \fig[height fraction]{name}: a figure from figures/, scaled to the frame; in the article it is a
% numbered figure whose caption is \figcap@<name> from figcaptions.tex.
% Captions of the figures in the article (the slides show none). One entry per figure file.
\expandafter\newcommand\csname figcap@fig-example\endcsname{A unit square cut into two triangles by a diagonal; the shaded triangle is the one used in the statement.}

\ifdeck
  \newcommand{\fig}[2][0.55]{\begin{center}\adjustbox{width=0.95\linewidth,height=#1\textheight,keepaspectratio}{\input{figures/#2}}\end{center}}
  \newenvironment{withfig}[1]{\gdef\wfname{#1}\par\begin{columns}[T,onlytextwidth]\column{0.64\textwidth}}{\column{0.34\textwidth}\vspace{0.3em}\centering\adjustbox{max width=\linewidth,max height=0.62\textheight}{\input{figures/\wfname}}\end{columns}}
\else
  \newcommand{\fig}[2][0.3]{\begin{figure}[ht]\centering\adjustbox{max width=0.9\linewidth,max height=#1\textheight}{\input{figures/#2}}\caption{\csname figcap@#2\endcsname}\label{fig:#2}\end{figure}}
  \newenvironment{withfig}[1]{\gdef\wfname{#1}}{\fig[0.3]{\wfname}}
\fi

% ---------- statement environments (same names as the paper, plus idea) ----------
\ifdeck
  \newtheorem{thm}{Theorem}
  \newtheorem{prop}[thm]{Proposition}
  \newtheorem{lem}[thm]{Lemma}
  \newtheorem{cor}[thm]{Corollary}
  \newtheorem{ideaenv}[thm]{Idea}
  \newtheorem{defn}[thm]{Definition}
  \newtheorem{ex}[thm]{Example}
  \newtheorem{rmk}[thm]{Remark}
  \newenvironment{idea}[1][]{\ifx\relax#1\relax\begin{ideaenv}\else\begin{ideaenv}[#1]\fi}{\end{ideaenv}}
\else
  \usepackage{amsthm}
  \newtheorem{thm}{Theorem}[section]
  \newtheorem{prop}[thm]{Proposition}
  \newtheorem{lem}[thm]{Lemma}
  \newtheorem{cor}[thm]{Corollary}
  \newtheorem{ideaenv}[thm]{Idea}
  \theoremstyle{definition}
  \newtheorem{defn}[thm]{Definition}
  \newtheorem{ex}[thm]{Example}
  \newtheorem{rmk}[thm]{Remark}
  \newenvironment{idea}[1][]{\ifx\relax#1\relax\begin{ideaenv}\else\begin{ideaenv}[#1]\fi}{\end{ideaenv}}
  \usepackage[margin=2.8cm]{geometry}
  \usepackage[hidelinks]{hyperref}
  \usepackage{microtype}
  \usepackage{booktabs}
  \usepackage{enumitem}
  \usepackage{caption}
  \captionsetup{font=small,labelfont=bf}
  \setlength{\parskip}{0.4em}
  \usepackage{cleveref}
  \crefname{thm}{Theorem}{Theorems}
  \crefname{prop}{Proposition}{Propositions}
  \crefname{lem}{Lemma}{Lemmas}
  \crefname{cor}{Corollary}{Corollaries}
  \crefname{ideaenv}{Idea}{Ideas}
  \crefname{defn}{Definition}{Definitions}
  \crefname{figure}{Figure}{Figures}
  \crefname{section}{Section}{Sections}
\fi

% ---------- deck look: one accent colour, grey blocks, serif mathematics ----------
\ifdeck
  \definecolor{accent}{HTML}{1F2A44}
  \usefonttheme[onlymath]{serif}
  \setbeamercolor{normal text}{fg=black}
  \setbeamercolor{structure}{fg=accent}
  \setbeamercolor{frametitle}{fg=accent}
  \setbeamerfont{frametitle}{size=\Large,series=\bfseries}
  \setbeamertemplate{frametitle}{\vskip1.2ex\usebeamerfont{frametitle}\usebeamercolor[fg]{frametitle}\insertframetitle\par\vskip0.4ex\hrule height 0.5pt\vskip0.5ex}
  \setbeamertemplate{blocks}[default]
  \setbeamercolor{block title}{fg=black,bg=black!7}
  \setbeamercolor{block body}{fg=black,bg=}
  \setbeamerfont{block title}{series=\bfseries}
  \setbeamertemplate{itemize item}{\raisebox{0.2ex}{\scriptsize$\bullet$}}
  \setbeamertemplate{itemize subitem}{\raisebox{0.2ex}{\scriptsize$\circ$}}
  \setbeamercolor{itemize item}{fg=accent}
  \setbeamercolor{enumerate item}{fg=accent}
  \setbeamertemplate{enumerate item}{\insertenumlabel.}
  \setlength{\leftmargini}{1.2em}
\fi
